\sectionkicker{Theme synthesis}
\subsection*{横断比較 — Live environmentを何で置き換えるか}
\addcontentsline{toc}{subsection}{横断比較 — Live environmentを何で置き換えるか}
DynaWebとSYNTHAGENTはどちらもreal environment依存を下げるが、代替しているものが異なる。world modelによるrolloutと、task・tool・user・rewardの合成を分けると設計上のtrade-offが見える。
\medskip
\begin{center}
\small
\renewcommand{\arraystretch}{1.16}
\begin{tabularx}{\linewidth}{@{}p{0.20\linewidth}X@{}}
\toprule
観察軸 & 1月の一次資料・論文から読めること \\
\midrule
world representation & DynaWebはAgent actionから次のpage representationを予測するweb world modelを学習する。live webを直接反復する代わりに、状態遷移を近似するlearned environmentを訓練対象へ持ち込む。 \cite{src-85bc3e4f7b39} \\
model-based rollout & 学習したworld modelはmodel-based reinforcement learningのrollout用synthetic environmentとして使われる。DynaWebの代替対象は主にweb transitionであり、task specification全体を合成する方式とは異なる。 \cite{src-85bc3e4f7b39} \\
DynaWeb evaluation & WebArenaとWebVoyagerでの改善は論文著者が報告する実験結果である。open-source web-agent model上の結果を、異なるlive serviceやproduction Agent全般の優位性へ拡張しない。 \cite{src-85bc3e4f7b39} \\
task / mock tools & SYNTHAGENTはtool-use taskとmock tool environmentを合成する。learned world transitionではなく、Agentが解くtaskと呼び出すtool側をsyntheticに構成する点がDynaWebとの大きな差である。 \cite{src-01d8ab4a2c6f} \\
user / reward & SYNTHAGENTは不足するprivate informationを補うLLM user simulatorとtask-level rubric rewardも構成する。environmentだけでなくuser interactionと評価信号までsynthetic stackに含める。 \cite{src-01d8ab4a2c6f} \\
SYNTHAGENT evaluation & 14のmath・search・tool-use datasetでの改善や小型modelが大型baselineを上回る例はauthor-reported resultである。synthetic trainingの一般的優位性を独立再現したものとしては扱わない。 \cite{src-01d8ab4a2c6f} \\
simulation gap & learned web world modelもsynthetic tool environmentも、dynamic・adversarialなlive systemの障害や攻撃条件を過小表現し得る。安定したtraining surfaceと実環境transferは別問題として残る。 \cite{src-85bc3e4f7b39,src-01d8ab4a2c6f} \\
chronology boundary & 両taskのlocatorは1月後に更新されたv2を含むため、Retrospectiveで後日読める本文を参照しても、月次chronologyはそれぞれのJanuary original submission dateへ固定する。 \cite{src-85bc3e4f7b39,src-01d8ab4a2c6f} \\
\bottomrule
\end{tabularx}
\renewcommand{\arraystretch}{1.0}
\end{center}
\normalsize
\begin{claimboundary}[この比較の境界]
この表は本号で既に検証・参照している一次資料と論文を横断比較の形へ再配置したものである。vendor / project / author が報告した性能値や優位性は独立再現済みの一般則へ変換せず、本文とSource Notesの帰属境界をそのまま維持する。
\end{claimboundary}
